\usepackage{fontspec}
\setmainfont{Times New Roman}
\setmonofont{Courier New}
\newfontfamily{\cyrillicfonttt}{Courier New}

\usepackage{polyglossia}
\setdefaultlanguage{russian}
\setotherlanguage{english}

% Поля согласно ГОСТ и методичке: левое 30 мм, правое 10 мм, верхнее и нижнее 20 мм
\usepackage[left=30mm, right=10mm, top=20mm, bottom=20mm]{geometry}

\usepackage{indentfirst} % Абзацный отступ
\setlength{\parindent}{1.25cm}

\usepackage{setspace} % Полуторный интервал
\onehalfspacing

\usepackage{titlesec} % Настройка заголовков
% Заголовки с абзацного отступа, строчными (начиная с прописной)
\titleformat{\chapter}[block]
  {\normalfont\bfseries\Large}{\thechapter}{1em}{}
\titlespacing*{\chapter}{\parindent}{*0}{15mm}

\titleformat{\section}[block]
  {\normalfont\bfseries\large}{\thesection}{1em}{}
\titlespacing*{\section}{\parindent}{8mm}{15mm}

\titleformat{\subsection}[block]
  {\normalfont\bfseries\normalsize}{\thesubsection}{1em}{}
\titlespacing*{\subsection}{\parindent}{8mm}{15mm}

% Настройка содержания
\usepackage{tocloft}
\renewcommand{\cfttoctitlefont}{\hfill\normalfont\bfseries\Large}
\renewcommand{\cftaftertoctitle}{\hfill\mbox{}}
\renewcommand{\cftchapleader}{\cftdotfill{\cftdotsep}}

% Графика и математика
\usepackage{graphicx}
\usepackage{amsmath, amssymb}

% Таблицы
\usepackage{array}
\setlength{\tabcolsep}{4pt}

% Списки
\usepackage{enumitem}
\setlist{nosep}
\setlist[itemize]{label=--} % дефис для маркированных списков
\setlist[enumerate]{label=\arabic*)} % 1), 2) для нумерованных

% Подписи к рисункам и таблицам
\usepackage{caption}
\captionsetup[table]{justification=raggedright, singlelinecheck=false, labelsep=endash, format=hang}
\captionsetup[figure]{justification=centering, labelsep=endash, name=Рисунок}

% Листинги кода
\usepackage{listings}
\renewcommand{\lstlistingname}{Листинг}
\captionsetup[lstlisting]{justification=raggedright, singlelinecheck=false, labelsep=endash, format=hang}
\lstset{
    basicstyle=\ttfamily\small,
    breaklines=true,
    showstringspaces=false,
    frame=none, % невидимые границы, как в методичке
    numbers=none, % нумерация строк (по умолчанию выключена, можно включить numbers=left)
    xleftmargin=1.25cm, % абзацный отступ
    tabsize=4
}

% Нумерация формул в пределах раздела
\numberwithin{equation}{chapter}

% Для ссылок
\usepackage[hidelinks]{hyperref}

% Нумерация страниц справа внизу
\usepackage{fancyhdr}
\pagestyle{fancy}
\fancyhf{} % очистка всех полей
\renewcommand{\headrulewidth}{0pt} % без линии в колонтитуле
\fancyfoot[R]{\thepage}

% Предотвращение выхода текста за рамки (Overfull \hbox)
\tolerance=1000
\emergencystretch=3em
